\documentclass[12pt]{report}

\usepackage[T1]{fontenc}
\usepackage[utf8]{inputenc}
\usepackage{graphicx}
\usepackage{amsmath,amssymb,amsfonts}
\usepackage{txfonts}
\usepackage{pdfpages}

\renewcommand{\chaptername}{Rozdział}
\renewcommand{\contentsname}{Spis treści}
\renewcommand{\figurename}{Rys.}
\renewcommand{\tablename}{Tab.}
\renewcommand{\listfigurename}{Spis rysunków}
\renewcommand{\listtablename}{Spis tabel}
\renewcommand{\bibname}{Bibliografia}

\pagestyle{headings}

\setlength{\textwidth}{14cm}
\setlength{\textheight}{20cm}

\newtheorem{definition}{Definicja}
\newtheorem{example}{Przykład}[chapter]
\newtheorem{corollary}{Wniosek}[chapter]

\begin{document}

\includepdf[pages={1}]{resources/title_page.pdf}

\tableofcontents

\chapter{Wstęp} \label{wstep}

\section{Uzasadnienie wyboru tematu}

Wraz z rozwojem informatyki proces tworzenia oprogramowania wymagał od informatyków co raz większej wiedzy w poszczególnych dziedzinach takich jak na przykład bazy 
danych czy sieci komputerowe. Ciężko jest być specjalistą w każdej z tych dziedzin, dlatego też podczas tworzenia oprogramowania programiści co raz częściej sięgają po
różnego rodzaju narzędzia programistyczne, które mają za zadanie ułatwić im pracę. 

Przykładem takich narzędzi są aplikacje szkieletowe wykorzystywane jako fundamenty dla tworzonych aplikacji czy też biblioteki udostępniające zestawy określonych funkcji. 
Oczywiście nie zawsze mamy możliwość skorzystania z wspominanych narzędzi, jednak jeśli taka istnieje warto wziąć to pod uwagę podczas fazy planowania tworzenia 
oprogramowania, ponieważ decydując się na korzystanie z nich jesteśmy w stanie zaoszczędzić sporo czasu, a także uniknąć wielu błędów związanych z nieznajomością danej
dziedziny.

Wraz z kolegą z tego samego roku studiów -- Sebastianem Florkiem, zdecydowaliśmy się w ramach pisamia pracy dyplomowej na wspólne stworzenie aplikacji szkieletowej, 
która będzie realizowała mapowanie obiektowo-relacyjne w języku C++. Wybór języka programowania związany był z dość dobrą jego znajomością, a także z pewnym
doświadczeniem w programowaniu w tym języku nabytym w trakcie studiów. Mapowanie obiektowo-relacyjne jest to obecnie zagadnienie bardzo powszechne, szczególnie w 
językach programowania takich jak C\# czy Java. Programiści C++ nie mają już tak dużego wyboru wśród dostępnych narzędzi służących do mapowania 
obiektowo-relacyjnego, co tłumaczy nasz wybór.

Podział naszej pracy zostanie opisany w dalszych rozdziałach, jednak na wstępie warto zaznaczyć, że tematem pracy Sebastiana jest generator opisu mapowania 
obiektowo-relacyjnego, który jest wstępnym etapem pracy naszej aplikacji szkieletowej. Celem wspomnianego generatora jest wygenerowanie pliku projektu, a także plików
nagłówkowych oraz klas na podstawie istniejącej już bazy danych. Moim zadaniem będzie samo mapowanie obiektowo-relacyjne, które będzie się opierało na wcześniej
wygenerowanych plikach. 

\section{Problematyka i zakres pracy}

Programowanie obiektowe jest obecnie jednym z najpopularniejszych paradygmatów programowania, a pojęcia takie jak klasa czy obiekt znane są wszystkim programistom.
Podobnie jest z relacyjnym modelem organizacji baz danych i terminami takimi jak relacja czy krotka. Chcąc wykorzystać oba te podejścia w jednej aplikacji musimy zadbać o 
obustronną konwersję pomiędzy danymi z tabel relacyjnej bazy danych a obiektami aplikacjii. Tym właśnie zajmuje się mapowanie obiektowo-relacyjne, które wraz z tworzeniem 
aplikacji szkieletowych w języku C++ jest główną problematyką niniejszej pracy.

W tym momencie powstaje pytanie czy na prawdę warto korzystać z bibliotek i aplikacji szkieletowych służących do mapowania obiektowo relacyjnego? Odpowiedź nie jest
jednoznaczna w wszystkich przypadkach, ale warto wymienić jego podstawowe wady i zalety, których głębsza analiza znajduje się w dalszej części pracy. Zacznijmy od zalet 
wykorzystania narzędzi ORM\footnote{mapowanie obiektowo-relacyjne (ang. Object-Relational Mapping)}:

\begin{itemize}
\item Znacznie zredukowana ilość pracy wymagana na oprogramowanie dostępu do bazy danych.
\item Nieobowiązkowa znajomość języka SQL\footnote{strukturalny język zapytań (ang. Structured Query Language)}. W celu tworzenia zapytań do bazy
danych korzystamy z interfejsu udostępnianego przez daną bibliotekę czy też aplikację szkieletową.
\item Uniezależnienie się od rodzaju systemu zarządzania bazą danych, możemy korzystać równie dobrze z MySQL, Oracle, PostgreSQL czy też Microsoft SQL Server
niekoniecznie posiadając o nich rozbudowaną wiedzę.
\item Aby skorzystać z transakcji, połączenia z bazą danych a także wieu innych funkcjonalności baz danych wystarczy zazwyczaj wywołać pojedyńczą metodę.
\item Cały model danych przechowywany jest w jednym miejscu oraz nie jest ściśle związany z wykorzystywanym systemem zarządzania bazą danych dzięki czemu łatwiej 
jest nam wprowadzać kolejne modyfikacje w kodzie oraz nawet zmieniać system zarządzania bazą danych na inny.
\end{itemize}

Wszędzie gdzie pojawiają się zalety mamy też do czynienia z wadami, nie inaczej jest w tym przypadku:

\begin{itemize}
\item Konfiguracja jest najczęściej skomplikowana i wymaga sporo czasu.
\item Podobnie jest z użytkowaniem, aby robić to w sposób optymalny musimy dobrze poznać daną bibliotekę czy też aplikację szkieletową.
\item Proste zapytania są obsługiwane bardzo sprawnie, jednak gdy przetwarzamy duże ilości złożonych zapytań wydajność nie dorówna nigdy zapytaniom napisanym
przez specjalistę znającego język SQL.
\item Abstrakcja wprowadzona przez narzędzia ORM może okazać się uciążliwa, ponieważ nie zawsze zdajemy sobie sprawę z tego co dzieje się za kulisami w trakcie
wykonywania poszczególnch operacji.
\end{itemize}

Naszym głównym celem jest uczynienie Qubica dobrą alternatywą dla nielicznych, ale istniejących już bibliotek oraz aplikacji szkieletowych realizujących mapowanie 
obiektowo-relacyjne w języku C++. Wszystkie nam obecnie znane rozwiązania są dostępne za darmo, jednak nie udostępniają one generatora opisu będącego w stanie 
wygenerować cały projekt aplikacji a ich interfejsy nie należą do intuicyjnych. Wprowadzenie generatora oraz intuicyjnego interfejsu użytkownika powinno uczynić Qubica
istotną alternatywą dla istniejących już narzędzi.

Analizę istniejących rozwiązań przedstawia kolejny rozdział a zestawienie z rezultatami naszej pracy znajduje się natomiast w końcowej części pracy, gdzie zostaną
przedstawione uzyskane przez nas wyniki oraz podsumowanie wykonanej przez nas pracy.

\section{Cele pracy} % Podrozdział około 1 strony, czyli 1800 znaków. Każdy cel opisany w 2-3 zdaniach, bez używania skrótów.

Do najważniejszych celów niniejszej pracy dyplomowej należą:

\begin{itemize}
\item Analiza istniejących bibliotek oraz aplikacji szkieletowych realizujących mapowanie obiektowo-relacyjne w języku C++.
\item Stworzenie własnej aplikacji szkieletowej realizującej mapowanie obiektowo-relacyjne w języku C++.
\item Analiza porównawcza szybkości działania przykładowej aplikacji stworzonej w oparciu o Qubic a o inne istniejące narzędzia.
\item Analiza porównawcza ilości kodu przykładowej aplikacji stworzonej w oparciu o Qubic a o inne istniejące narzędzia.
\end{itemize}

Do celów części praktycznej należą:

\begin{itemize}
\item Stworzenie intuicyjnego interfejsu użytkownika -- aby sprawdzić w jakim stopniu udało się zrealizować założenia zostanie przeprowadzony test polegający na napisaniu
tej samej aplikacji wykorzystującej Qubica oraz inne aplikacje szkieletowe i bibioteki, a następnie porównaniu ilości linii kodu stworzonych aplikacji.
\item Stworzenie generatora opisu mapowania obiektowo-relacyjnego -- jest to przedmiotem pracy Sebastiana, naszym wspólnym celem jest integracja obu modułów.
\item Poprawne zrealizowanie założeń mapowania obiektowo-relacyjnego, a także jak najlepsza optymalizacja zapytań -- aby sprawdzić w jakim stopniu udało się zrealizować
założenia zostanie przeprowadzony test polegający na napisaniu tej samej aplikacji wykorzystującej Qubica oraz inne aplikacje szkieletowe i bibioteki, a następnie porównaniu
ich wydajności.
\item Uczynienie konfiguracji Qubica jak najprostszą.
\end{itemize}

\section{Metoda badawcza}

\subsubsection{Studia literaturowe} % Podrozdział około 1 strony, czyli 1800 znaków.Opis w 2-3 zdaniach. Wymagane odniesienia do bibliografii!

Podstawowa literatura wykorzystana podczas pisania niniejszej pracy dotyczy specyfikacji języka C++ oraz zapytań w języku SQL, są to powszechne zagadnienia, więc wybór
pozycji książkowych jest spory. Materiały dotyczące samego mapowania obiektowo-relacyjnego pochodzą najczęściej z źródeł elektronicznych w języku angielskim, co jest
związane z faktem, że wiedza ta wciąż się rozwija.

\subsubsection{Analiza istniejących rozwiązań}

\subsubsection{Stworzenie własnej aplikacji szkieletowej}

Praktyka jest najczęściej najlepszą z dostępnych metod nauki i to właśnie podczas tworzenia własnej aplikacji szkieletowej realizującej mapowanie obiektowo-relacyjne można
się najbardziej z danym tematem zapoznać. Wszystkie problemy, które pojawiały się poczas pisania Qubica musiały zostać w pewien sposób rozwiązane i to właśnie analiza
tych problemów i ich rozwiązywanie było główną metodą badawczą wykorzystaną podczas pisania niniejszej pracy.

\subsubsection{Analiza porównawcza oraz testy}

\section{Przegląd literatury w dziedzinie}

\subsubsection{Literatura dotycząca języka C++}

\subsubsection{Literatura dotycząca języka SQL}

\subsubsection{Literatura dotycząca mapowania obiektowo-relacyjnego}

\subsubsection{Literatura dotycząca aplikacji szkieletowej Qt}

\section{Układ pracy} % Cały podrozdział około 1 strony.

\chapter{Zagadnienia teoretyczne dotyczące mapowania obiektowo-relacyjnego oraz aplikacji szkieletowych} \label{teoria}

\section{Podstawowe definicje} % Dokładny opis terminologii  pojęć zasadniczych dla tematu pracy, którymi autor będzie się posługiwał przy realizacji głównych celów pracy.

\chapter{Technologie i metody użyte w części badawczej} \label{technologie}

\chapter{Aplikacja szkieletowa Qubic} \label{qubic}

\section{Dokumentacja techniczna}

\section{Dokumentacja użytkownika}

\section{Przykładowa aplikacja wykorzystująca Qubica}

\section{Testy oraz ich wyniki}

\section{Perspektywy rozwoju Qubica}

W celu dalszego rozwoju stworzonej aplikacji szkieletowej warto rozważyć wprowadzenie następujących usprawnień:

\begin{itemize}
\item System adnotacji -- obecnie na opis mapowania składają się odpowienie nazwy funkcji oraz makra Qt. Istnieje jednak możliwość wprowadzenia własnych makr, które
miałyby opisywać mapowanie pomiędzy nazwami tabel z baz danych a odpowiednimi polami klas napisanych z języku C++. Dzięki temu zaistniałaby możliwość uniezależnienia
nazw pól klas od nazw tabel w bazie danych.
\item Interfejs zapytań -- choć jest już zaimplementowany, nadal nie udostępnia on wszystkich możliwych funkcjonalności języka SQL. Implementacja obsługi takich poleceń
jak JOIN czy UNION z pewnością byłaby dodatkowym atutem.
\item Identyfikacja tabel także za pomocą kluczy złożonych -- w tej chwili tabele identyfikowane są za pomocą kluczów głównych, co z kolei wymusza ich nadawanie w każdej
z tabel.
\item Pamięć podręczna -- wprowadzenie pamięci podręcznej może znacznie polepszyć wydajność w przypadku ciągłych operacji na tych samych danych.
\item Konfiguracja z poziomu kodu -- obecnie większość konfiguracji jest zapisana w plikach konfiguracyjnych i tylko tam może być zmieniana, w celu rozwoju wprowadzenie
dodatkowej możliwości jego konfiguracji wydaje się być dobrym pomysłem.
\item Wsparcie dla różnych rodzajów baz danych -- wprowadzenie tego usprawnienia ogranicza się do implementacji kilku interfejsów dla innych niż MySQL rodzajów baz 
danych. Biorąc pod uwagę możliwość wzorowania się na zaimplementowanej już logice nie powinno to stworzyć problemu gdy zaistnieje taka konieczność.
\item Serializacja danych -- dodanie możliwości serializacji może okazać się użyteczne w przypadku pracy z dużymi ilościami danych, w tym celu można skorzystać z wielu
istniejących już bibliotek udostępniających tę możliwość.
\item Internacjonalizacja -- w tej chwili wszystkie logi zlokalizowane są w języku angielskim, istnieje jednak możliwość zmiany obecnego stanu poprzez wykorzystanie modułu
translacji udostępnianego przez Qt.
\item Wielowątkowość -- wykorzystanie wielowątkowości w przypadku mapowania-obiektowo relacyjnego z pewnością nie należy do najłatwiejszych zadań, jednak znacznie
może to usprawnić wykonywanie bardziej wymagających operacji.
\end{itemize}

\chapter{Podsumowanie} \label{podsumowanie}

\section{Dyskusja wyników}

\section{Perspektywy rozwoju pracy} % Jak można kontynuować tę pracę, zwłaszcza pod kątem magisterskich i dalej? Co jeszcze powinno być zrobione lub ulepszone?
% Co należy zmienić lub poprawić w pracy z dzisiejszego punktu widzenia?

%%%%%%%%%%%%%%%%%%%%%%%%%%%%%%%%%%%%%%%%%%%%%%%%%%%%%%%%%%%%%%%%%%%%%%%%%%%%
%%%%%%%%%%%%%%%%%%%%%%%%%%%%%%%%%%%%%%%%%%%%%%%%%%%%%%%%%%%%%%%%%%%%%%%%%%%%
%%%%%%%%%%%%%%%%%%%%%%%%%%%%%%%%%%%%%%%%%%%%%%%%%%%%%%%%%%%%%%%%%%%%%%%%%%%%

\newpage

Tematem pracy jest: ....., zaś za główny cel przyjęto ...... . \\
Rozdział \label{rozdz.wstep} zawiera wstęp i cele pracy. W rozdziale drugim
opisano/...... w Rozdziale 3. zawarto............ Rozdział 4. przedstawia..... \\
W podsumowaniu pracy przedstawiono..........................., z czego wynika,
że ................  \\
Najważniejszym wnioskiem/wynikiem/rezultatem pracy jest..................\\ {\bf wyraźnie określić
CO TO JEST}. \\

%% http://webhosting.pl/Czy.ORM_y.odejda.wkrotce.do.lamusa
%% http://zeszyty-naukowe.wwsi.edu.pl/zeszyty/zeszyt4/Mapowanie_Obiektowo-Relacyjne_ORM_-_Czy_Tylko_Dobra_Idea.pdf
%% http://sens.e-informatyka.pl/wp-content/uploads/WIRP2/Hibernate.pdf

Mapowanie obiektowo-relacyjne (ang. Object-Relational Mapping ORM) to sposób odwzorowania obiektowej architektury systemu informatycznego na bazę danych (lub inny element systemu) o relacyjnym charakterze. Implementacja takiego odwzorowania stosowana jest m.in. w przypadku, gdy tworzony system oparty jest na podejściu obiektowym, a system bazodanowy (System Zarządzania Bazą Danych) operuje na relacjach. Z ORM związany jest szereg problemów wydajnościowych.
ORM czyli Object-Relational Mapping jest to rozwiązanie w którym dane są mapowane i zwracane w postaci obiektów.

Odwzorowania obiektowo relacyjne
A) ORM pozwala na reprezentację danych pobieranych z bazy w postaci obiektów (kolekcji obiektów, obiektów powiązanych przez referencje). Wiązanie to jest zdefiniowane w sposób deklaratywny (adnotacje lub XML). Izoluje to kod aplikacji (opierający się na paradygmacie obiektowym) od bazy danych (opierającej się na paradygmacie relacyjnym).
B) ORM zapewnia obiektom trwałość (ang. persistence). Modyfikacje struktury obiektowej są zapisywane w bazie.
A) I B) pozwala na wysokopoziomowy dostęp do bazy. PROBLEM: Niezgodność modeli (impedance mismatch) relacje w bazie są reprezentowane inaczej niż relacje między obiektami.  Ponadto baza nie pozwala na dziedziczenie a obiekty tak. ORM wprowadzają sposoby radzenia sobie z tą niezgodnością.

ORM - projektowanie aplikacji
Tworzenie warstwy trwałości korzystającej z ORM może być realizowane na 3 sposoby:
A) Mapowanie do przodu (forward mapping) - mamy klasy i na ich podstawie będziemy tworzyć bazę danych. To jest najprostsze i nie wymaga zbyt dużej ilości metadanych w kodzie (domyślnie nazwa właściwości jest mapowana do nazwy pola)
B) Mapowanie do tyłu (reverse mapping) - mamy bazę danych i na jej podstawie będziemy tworzyć klasy. To jest trudniejsze - wybaga definiowania większej ilości metadanych.
C) Mapowanie Meet in The Middle - mamy i schemat bazy i klasy i chcemy je do siebie wzajemnie dopasować.

http://stackoverflow.com/questions/1279613/what-is-an-orm-and-where-can-i-learn-more-about-it

\section{Istniejące rozwiązania w dziedzinie}

W tym podrozdziale zostaną opisane.....
\subsection{Sprzęt}
.........................
\subsection{Oprogramowanie i wdrożone systemy}
.....................................
\subsection{}
...................

\section{Wady i słabe punkty istniejących rozwiązań}
\subsection{Efektywność}
..........................
\subsection{Utrudniony dostęp}
..............
\subsection{Wysokie koszty}
...............................

\chapter{Dalsze uwagi o edycji i~formatowaniu pracy}
Pracę w \LaTeX'u najlepiej składać w szablonie {\tt report}, ze względu na jendostronny wydruk (jak w {\tt article}) i możliwość dzielenia pracy na rozdziały, a co za tym idzie, tworzenia spisu treści, spisu tabel, rysunków. 
\begin{example}
Przyklad
\end{example}

\begin{corollary}
Wniosek
\end{corollary}
\section{Bibliografia i przypisy}
Spis litertury dołącza się  w \LaTeX'u automatycznie na końcu pracy (zob.
komenda {\tt begin{thebibliography}}). Informacje o sposobie cytowania zawarte
są na stronie Bibilioteki Głównej PŁ\\ także udostępnione na
\underline{\tt http://ics.p.lodz.pl/\textasciitilde aniewiadomski}. \\

\indent Przykład cytowania --  jak podaje praca \cite{kacprzyk86}, ......,
jednakże autorzy [2] twierdzą, iż.....\\


\indent Za każdym razem, kiedy w pracy pojawia się treść na podstawie jakiegoś
tekstu źródłowego czyjegoś autorstwa, oznaczamy takie miejsce
przypisem\footnote{Treść przypisu pierwszego}. Przypis zawierać musi  numer
jakim w spisie literatury, czyli bibliografii, oznaczono tę pracę, np.
tak\footnote{[3], ss. 3--6 (czyli praca trzecia w spisie literatury,
wykorzystany fragment znajduje się na stronach od 3. do 6.)}. {\bf
Wszystkie źródła tekstów, rysunków, danych, wykresów, schematów, kodów i
informacji wykorzystanych w pracy muszą być zamieszczone w bibliografii.
Wszystkie pozycje literatury zamieszczone w bibliografii muszą być cytowane w
treści pracy, na dowód, iż zostały rzeczywiście użyte przy pisaniu pracy.}

\subsubsection{Źródła elektroniczne}
Źródła elektroniczne, zwłaszcza internetowe należy cytować z należytą uwagą na
ich jakość. Nie cytujemy źródeł wątpliwej jakości lub wtórnie przekazujących
czy też powielających wiedzę zawartą w innych źródłach, np. fora internetowe lub
wikipedia.\\
\indent {\bf Wszystkie wykorzystane źródła elektroniczne powinny być przez Autora
pracy skopiowane \underline{w dniu ich wykorzystania} i
dołączone np. na CD/DVD do wersji drukowanej pracy.\\
\indent Odnośniki do źródeł elektronicznych muszą zawierać pełną ścieżkę, np. do
pliku lub rysunku, a nie jedynie domenowy adres portalu, np.\newline {\tt
http://serwer.com/temat/podtemat/katalog/plik\_strony.html} (stan na dzień:
2009-12-05)\\
ale nie\\
{\tt www.portal.pl}.} (!!!!!)

Niedochowanie tego wymogu może stać się powodem odrzucenia pracy ze wzglę\-dów
formalnych (,,brak możliwości weryfikacji źródeł wykorzystanych w pracy''). 

\section{Polskie akapity, cudzysłowy, itp.}
Akapity stosujemy zawsze z wcięciem, ale bez wiersza odstępu pomiędzy akapitami. 
Ta forma jest przyjęta dla publikacji polskojęzycznych. {\bf W szczególnych
przypadkach (także w tym szablonie)} akapit występujący bezpośrednio po tytule
rozdziału, sekcji, podsekcji itp. NIE JEST WCIĘTY.\\
\indent Ten akapit JEST WCIĘTY. NIE MA także PUSTEGO WIERSZA pomiędzy tym
akapitem a poprzednim. \\
\indent Podobne uwagi dotyczą wszystkich innych elementów formatowania pracy --
muszą być zgodne ze zwyczajami przyjętymi W JĘZYKU POLSKIM.	Np. cu\-dzysłowy
wyglądają tak: ,,cudzysłów'', ale nie ''cudzysłów'', albo też `cudzysłów' czy
``cudzysłów''. 



\section{Definicje i wyrażenia matematyczne}
\begin{definition} \label{def.definicja1}
Niech $\cal X$ będzie przestrzenią.....
\end{definition}

Do definicji odnieść sie można poprzez jej etykietę: jak podano w Def.~\ref{def.definicja1}

Przykładowe podkreślenie... \underline{tekst podkreślony}, pogrubienie: {\bf
tekst pogrubiony} oraz wyrożnienie {\em tekst wyróżniony, czyli kursywa}. 
Dalszy tekst rozdziału
Dalszy tekst rozdziału
Dalszy tekst rozdziału
Dalszy tekst rozdziału
Dalszy tekst rozdziału
Dalszy tekst rozdziału
Dalszy tekst rozdziału
Dalszy tekst rozdziału
Dalszy tekst rozdziału a teraz koniec linii... \\
\indent ... i nowy akapit. Akapity muszą być standardowo wcięte.  


Przykład wzoru matematycznego numerowanego

\begin{equation} \label{wzoreinsteina}
E=m\cdot c^2
\end{equation}

{\bf Wszystkie symbole matematyczne występujące w tekście ,,na bieżąco'',
czyli nieoznaczone numerem równania TAKŻE PISZEMY W TRYBIE MATEMA\-TYCZNYM, CZYLI
K U R S Y W Ą} : $a=b\cdot c$, ale nie: a= b*c (!!)\\

\indent Numeracja wzoru -- ZAWSZE w POSTACI (\#.\#\#)
Jak podaje wzor (\ref{wzoreinsteina}).... (koniec linii). \\
\indent Wyrażenia matematyczne można też wpisywać w wierszu -- używamy wów\-czas znaku '\$', który rozpoczyna i kończy wyrażenie, np. wg Einsteina $E=m\cdot c^2$...


\section{Jak wstawiać rysunki? tabele? }
A teraz pora na rysunek:
\begin{figure}[!t]
\centering
%\includegraphics[width=7cm]{1figmftall} 
\caption{Funkcja przynależności zbioru rozmytego -- Podpis ZAWSZE POD rysunkiem,
numeracja w postaci \#.\#\#. } (wypada podać źródło, czyli literaturę,
z której rysunek pochodzi, ewentualnie {\em opracowanie własne}.)
\label{fig.funkcja.przyn}
\end{figure}


\begin{table}[!t]
\centering
\caption{Tytuł tabeli ZAWSZE NAD TABELĄ, numeracja w formie \#.\#\#. (wypada podać źródło, czyli literaturę,
z której tabela pochodzi, ewentualnie {\em opracowanie własne}.)} 

\label{tabls1}

{\footnotesize 
\vspace{5mm}
\begin{tabular}{c c c c c}
\hline\noalign{\smallskip}
{\bf Alg.} & {\bf tytuł kolumny 1} & {\bf tytuł kolumny 1} & {\bf Tytuł kolumny
3} & {\bf ....}     \\

\hline\noalign{\smallskip}
a & b & c & d & e  \vspace{3mm} \\ 
\noalign{\smallskip}
 a & b & c & d & e \\

\noalign{\smallskip}
%%%
\end{tabular}
}
\end{table}


Rysunki i tabele nie powinny przekraczać 0.9 szerokości tekstu i zasadniczo
powinny występować na górze strony. \\

\indent Odnosić się do rysunku można poprzez jego etykietę ''label'', np. jak widać na rys. \ref{fig.funkcja.przyn}......

Jak widać, rysunek nie wypada w dokumencie w tym samym miejscu co w kodzie, choć czasem się tak zdarza. Jeśli potrzebujesz przenieść rysunek, zajrzyj do rozdzialu 2.11. manuala pt. {\em Wstawki}. 

\section{Listy wypunktowana i numerowana}

\begin{itemize}
\item pierwszy element listy wypunktowanej
\item drugi...
\item trzeci...
\end{itemize}


Nowy akapit z lista numerowaną. 
\begin{enumerate}
 \item pierwszy element listy NUMEROWANEJ
 \item drugi...
 \item trzeci...
 \item trzeci...
 \item trzeci...
 \end{enumerate}

\section{Przenoszenie wyrazów}
Skorzystaj z polecenia {\tt hyphenation}\\ w preambule dokumentu, lub dziel
wyrazy ,,ręcznie'' czyli właśnie tak jak tu: po\-dzie\-lo\-ne wy\-ra\-zy. 



{\em Tytuł tego rozdziału ma dwie wersje: zwykłą, (w kodzie: w nawiasach
klamrowych), która
pokazuje sie na stronie rozpoczynającej rozdział, oraz krótką (w kodzie: w nawiasach
kwadratowych), która pokazuje sie w spisie treści i w nagłówku}

W rozdziale \ref{etykietarozdzialu2} podano podstawy teoretyczne i ogólny zakres
pracy. W niniejszym rozdziale opisana zostanie technologia XYZ oraz metoda ABC
użyta w części praktycznej, patrz rozdział~\ref{rozdz.czesc.prakt}. 

\section{Sprzęt}
...................
\subsection{Element 1}
.........................
\subsection{Element 2}
......................

\section{Oprogramowanie}
..........................
\subsection{Serwer baz danych}
........................
\subsection{Środowisko zintegrowane}
..........................
\subsection{Oprogramowanie klienckie}

\section{Technologie i metodologie programistyczne}
..................
\subsection{Język programowania}
......................
\subsection{Biblioteki}
.......................
\subsection{Wzorce projektowe}
.......................

\section{Inne, np. narzędzia i metody symulacji, }

\chapter{Biblioteka programistyczna Qubic} \label{rozdz.czesc.prakt}
Ta część pracy może być podzielona na więcej rozdziałów, np kiedy autor chce
w~szczególności podkreślić któryś z etapów projektu. W zależności od tematu i~celów pracy, pewne sekcje można dodać (np. przy projektowaniu sieci, instalacji
i~konfiguracji serwerów usług sieciowych), inne zaś pominąć.

\section{Analiza wymagań}
\subsection{Studium możliwości}
\subsection{Wymagania funkcjonalne}
.................
\subsection{Ograniczenia projektu}

\section{Projekt}
\subsection{Projekt warstwy danych}

\begin{enumerate}
\item normalizacje baz danych
\item projekt bazy/baz 
\item grupy użytkowników i ich prawa dostępu do danych (zależne od implementacji bazy)
\item ew. diagramy klas warstwy danych
\end{enumerate}
\subsection{Projekt warstwy logiki}
\begin{enumerate}
\item Diagramy i scenariusze przypadków użycia
\item Diagramy przepływu danych (lub ich odpowiedniki)
\item ew. diagramy klas, wzorce projektowe itp.
\end{enumerate}

\subsection{Projekt warstwy interfejsu użytkownika}
\subsubsection{Wybór środowiska i platformy działania}
\subsubsection{Rodzaj aplikacji (klient-serwer, thick/thin client, aplikacja
,,biurkowa'', usługa, klient hybrydowy, itp.}
\subsubsection{Technologie projektowania i realizacji interfejsu użytkownika,
np. biblioteki}


\section{Implementacja: punkty kluczowe}

\section{Testy i wdrożenie}
\subsection{Testy wydajności}
\subsection{Testy regresyjne}
\subsection{Testy bezpieczeństwa}
\subsection{Dalsze testy}
\subsection{Testy...}

\section{Konserwacja i inżynieria wtórna}
Jak przebiega eksploatacja systemu/projektu? Jakie wady i zalety ujawniły się po
np. 2-miesięcznym okresie testowania i użytkowania? \\
\indent Jak można skorzystać z tej wiedzy praktycznej pod kątem roz\-bu\-do\-wy pracy? Jakie elementy systemu powinny zostać w pierwszej kolejności zmodyfikowane?  

\chapter{Podsumowanie}
\section{Dyskusja wyników}
Dzięki zrealizowaniu pracy poprawie uległa wydajność ....... Ponadto, o ?? \%
skrócony został czas ........, a koszty osiągnięcia zamierzonego efektu zostały
zmniejszone z ???pln do ???pln za godzinę/ dzień/ jednostkę sprzętu.........\\
\indent Które cele pracy udało sie zrealizować? co z tego wynika? Które cele
pracy pozostały niezrealizowane i dlaczego? 

\section[Ocena możliwości wdrożenia...]{Ocena możliwości wdrożenia proponowanych
\newline rozwiązań...}
... ich wartość praktyczna, lokalne i globalne możliwości zastosowania, kwestia
praw autorskich do powstałych produktów, itp. 


\addcontentsline{toc}{chapter}{Bibliografia} 
\begin{thebibliography}{99}
\bibitem{kacprzyk86}
Kacprzyk J. (1986) Fuzzy sets in system analysis.  PWN, Warsaw (in Polish). % http://www.ipipan.waw.pl/~subieta/artykuly/slownik_obiektowosci/hasla_slownika.html
\bibitem{kacprzyk99b}
Kacprzyk J., Strykowski P. (1999) Linguistic Data Summaries for Intelligent Decision Support, Proceedings of EFDAN'99. 4-th European Workshop on Fuzzy Decision Analysis and Recognition Technology for Management, Planning and Optimization, Dortmund, 1999, 3--12.
\bibitem{kacprzyk01d}
Kacprzyk J., Yager R. R. (2001) Linguistic summaries of data using fuzzy logic. International Journal of General Systems 30:133--154 

\end{thebibliography}

\addcontentsline{toc}{chapter}{Spis rysunków} 
\listoffigures

\addcontentsline{toc}{chapter}{Spis tabel} 
\listoftables


\addcontentsline{toc}{chapter}{Załączniki} 
\chapter*{Załączniki}
\begin{enumerate}
\item Załącznik nr 1
\item Załącznik nr 2
\item Załącznik nr 3
\end{enumerate}

%%%%%%%%%%%%%%%%%%%%%%%%%%%%%%%%%%%%%%%%%%%%%%%%%%%%%%%%%%%%%%%%%%%%%%%%%%%%
%%%%%%%%%%%%%%%%%%%%%%%%%%%%%%%%%%%%%%%%%%%%%%%%%%%%%%%%%%%%%%%%%%%%%%%%%%%%
%%%%%%%%%%%%%%%%%%%%%%%%%%%%%%%%%%%%%%%%%%%%%%%%%%%%%%%%%%%%%%%%%%%%%%%%%%%%

\end{document}